\documentclass[UTF8,aspectratio=169]{ctexbeamer}
\usetheme{Madrid}
\usecolortheme{default}
\usepackage{amsmath}
\usepackage{booktabs}
\usepackage{siunitx}
\usepackage{graphicx}
\usepackage{array}
\usepackage{xcolor}
\usepackage{xurl}
\usepackage{colortbl}
\usepackage{bm}

\newcommand{\good}[1]{\textcolor{green!45!black}{#1}}
\newcommand{\warn}[1]{\textcolor{red!65!black}{#1}}
\newcommand{\bluenote}[1]{\textcolor{blue!65!black}{#1}}
\newcommand{\src}[1]{\vfill{\tiny\textcolor{gray!75!black}{Source: \url{#1}}}}
\newcolumntype{L}[1]{>{\raggedright\arraybackslash}p{#1}}
\newcolumntype{R}[1]{>{\raggedleft\arraybackslash}p{#1}}
\newcolumntype{C}[1]{>{\centering\arraybackslash}p{#1}}

\definecolor{winrow}{RGB}{220,240,220}
\definecolor{warnrow}{RGB}{255,235,230}

\title[Stage-1 Feature Ablation]{Stage-1 Feature Ablation\\LONO Validation of Physics-Based Scaling}
\author{Mie Spray Postprocessing Project}
\date{\today}

\begin{document}

\begin{frame}
  \titlepage
\end{frame}

% ──────────────────────────────────────────────────────────────
\begin{frame}{Message}
\small
\textbf{Claim:} replacing raw pressure and diameter inputs with a single physics-based amplitude scale $A$ absorbs the dominant source of extrapolation error and produces a strictly better Stage-1 model under leave-one-nozzle-out (LONO) evaluation.

\vspace{0.5em}
\textbf{This deck:}
\begin{itemize}
  \item recaps how $A_{\mathrm{scale}}$ encodes the Hiroyasu--Arai asymptote and why the $\Delta P^{0.5}$ exponent is preferred
  \item defines all nine feature variants tested in the ablation
  \item shows per-fold and mean Stage-1 physical MAE for all variants under 5-fold LONO
  \item confirms that (i) $\Delta P^{0.5}$ beats $\Delta P^{0.25}$, (ii) adding pressure residuals gives a small further gain, and (iii) re-introducing $d_n$ as a standalone input consistently \warn{hurts}
\end{itemize}

\vspace{0.5em}
\textbf{Validated pipeline:} \texttt{run\_stage1\_lono\_ablation.py}, seed 42, 5 folds
($n_{\mathrm{fold}} \in \{1038, 195, 81, 81, 81\}$ sample groups).
\end{frame}

% ──────────────────────────────────────────────────────────────
\begin{frame}{Background: The Amplitude Scale $A_{\mathrm{scale}}$}
\small
The Hiroyasu--Arai model predicts spray tip penetration as
\[
  S(t) \;\approx\; A \cdot t^{1/4},
  \qquad
  A \;=\; \Delta P^{\,\alpha}\,\rho_a^{-1/4}\,d_n^{1/2},
\]
where $\alpha = 1/4$ is the classical H-A exponent.
Data-driven regression on our dataset supports $\alpha = 1/2$ (the transitional-regime exponent), giving
\[
  \boxed{A_{\mathrm{scale}} = \Delta P^{0.5}\,\rho_a^{-0.25}\,d_n^{0.5}.}
\]

\vspace{0.4em}
\textbf{Why this matters for feature engineering:}
\begin{itemize}
  \item $d_n$ and $\Delta P$ are \emph{already encoded} in $A_{\mathrm{scale}}$; the MLP sees the
        amplitude-normalised residual $\hat{y} = S/A_{\mathrm{scale}}$
  \item the network only needs to learn the \emph{shape} variation not captured by $A$: spray
        angle, number of plumes, injection duration, back-pressure, and optionally
        log-pressure corrections
  \item re-adding $d_n$ or pressures as standalone features creates a redundant pathway that
        the network can exploit to \emph{memorise} nozzle identities rather than generalise
\end{itemize}
\end{frame}

% ──────────────────────────────────────────────────────────────
\begin{frame}{Variant Design: What Each Configuration Tests}
\small
\begin{center}
\renewcommand{\arraystretch}{1.18}
\begin{tabular}{@{}L{3.6cm} C{1.0cm} C{1.0cm} C{1.0cm} C{1.0cm} L{3.8cm}@{}}
\toprule
\textbf{Variant} & $\bm{\alpha}$ & \textbf{Pressure} & \textbf{Diam.} & \textbf{\#feat} & \textbf{What it tests} \\
\midrule
\texttt{legacy\_9\_no\_scale} & --- & direct & direct & 9 & v1 baseline; all raw features, no $A$-scaling \\
\addlinespace[0.3em]
\texttt{a\_dp025} & 0.25 & --- & --- & 5 & pure H-A scale; pressures and $d_n$ fully absorbed \\
\texttt{a\_dp025\_plus\_pressures} & 0.25 & log & --- & 7 & H-A scale + log-pressure residuals \\
\texttt{a\_dp025\_plus\_diameter} & 0.25 & --- & $z$ & 6 & H-A scale + explicit $d_n$ (redundant test) \\
\texttt{a\_dp025\_plus\_pressures\_diameter} & 0.25 & log & $z$ & 8 & H-A scale + both residuals \\
\addlinespace[0.3em]
\texttt{a\_dp050} & 0.50 & --- & --- & 5 & transitional scale; pressures and $d_n$ absorbed \\
\texttt{a\_dp050\_plus\_pressures} & \good{0.50} & \good{log} & --- & \good{7} & \good{production variant (winner)} \\
\texttt{a\_dp050\_plus\_diameter} & 0.50 & --- & $z$ & 6 & transitional scale + explicit $d_n$ (redundant) \\
\texttt{a\_dp050\_plus\_pressures\_diameter} & 0.50 & log & $z$ & 8 & transitional scale + both residuals \\
\bottomrule
\end{tabular}
\end{center}

\vspace{0.3em}
\footnotesize
Base features present in \emph{all} $A$-scale variants: normalised time, tilt angle $\theta$, plume count $n_p$, injection duration $t_i$, back-pressure $P_{\mathrm{cb}}$.
Log-pressure residuals: $\log P_{\mathrm{inj}}$, $\log P_{\mathrm{ch}}$ (z-scored).
\end{frame}

% ──────────────────────────────────────────────────────────────
\begin{frame}{Feature Columns: Explicit Enumeration}
\small
\vspace{0.2em}
\begin{columns}[T]
\column{0.48\textwidth}
\textbf{Base set} (all $A$-scale variants):
\begin{itemize}\setlength\itemsep{2pt}
  \item \texttt{time\_norm\_0\_5ms} --- normalised time $t$
  \item \texttt{tilt\_angle\_radian\_z} --- spray half-angle $\theta$
  \item \texttt{plumes\_z} --- hole count $n_p$
  \item \texttt{injection\_duration\_us\_z} --- $t_i$
  \item \texttt{control\_backpressure\_bar\_z} --- $P_{\mathrm{cb}}$
\end{itemize}

\vspace{0.5em}
\textbf{Optional add-ons:}
\begin{itemize}\setlength\itemsep{2pt}
  \item \texttt{log\_injection\_pressure\_bar\_z} --- $\log P_{\mathrm{inj}}$
  \item \texttt{log\_chamber\_pressure\_bar\_z} --- $\log P_{\mathrm{ch}}$
  \item \texttt{diameter\_mm\_z} --- $d_n$
\end{itemize}

\column{0.48\textwidth}
\textbf{Legacy 9 (no $A$-scaling):}
\begin{itemize}\setlength\itemsep{2pt}
  \item \texttt{time\_norm\_0\_5ms}
  \item \texttt{tilt\_angle\_radian\_z}
  \item \texttt{plumes\_z}
  \item \texttt{diameter\_mm\_z}
  \item \texttt{injection\_duration\_us\_z}
  \item \texttt{log\_injection\_pressure\_bar\_z}
  \item \texttt{log\_chamber\_pressure\_bar\_z}
  \item \texttt{log\_delta\_pressure\_bar\_z}
  \item \texttt{control\_backpressure\_bar\_z}
\end{itemize}
\end{columns}

\vspace{0.4em}
\footnotesize
\textbf{Key difference:} in $A$-scale variants the target is $S / A_{\mathrm{scale}}$ (dimensionless);
in legacy the target is $S$ directly (mm). The network in legacy mode must therefore learn
$\Delta P$ and $d_n$ sensitivity from scratch without the physics prior.
\end{frame}

% ──────────────────────────────────────────────────────────────
\begin{frame}{LONO Results: Summary (Mean Across 5 Folds)}
\small
\begin{center}
\renewcommand{\arraystretch}{1.22}
\begin{tabular}{@{}rlrr@{}}
\toprule
\textbf{Rank} & \textbf{Variant} & \textbf{MAE mean (mm)} & \textbf{MAE std (mm)} \\
\midrule
\rowcolor{winrow}
1 & \texttt{a\_dp050\_plus\_pressures} & \textbf{7.46} & 6.52 \\
2 & \texttt{a\_dp025\_plus\_pressures} & 8.22 & 8.55 \\
3 & \texttt{a\_dp050} & 8.27 & 6.08 \\
4 & \texttt{a\_dp025} & 9.83 & 7.59 \\
\midrule
\rowcolor{warnrow}
5 & \texttt{a\_dp050\_plus\_diameter} & 12.06 & 7.67 \\
\rowcolor{warnrow}
6 & \texttt{a\_dp025\_plus\_diameter} & 13.41 & 7.60 \\
\rowcolor{warnrow}
7 & \texttt{a\_dp050\_plus\_pressures\_diameter} & 14.26 & 9.69 \\
\rowcolor{warnrow}
8 & \texttt{legacy\_9\_no\_scale} & 14.42 & 9.43 \\
\rowcolor{warnrow}
9 & \texttt{a\_dp025\_plus\_pressures\_diameter} & 14.48 & 10.21 \\
\bottomrule
\end{tabular}
\end{center}

\vspace{0.3em}
\small
\good{Green}: LONO winner (ranks 1--4, all $A$-scale without $d_n$ redundancy). \quad
\warn{Red}: variants that re-introduce $d_n$ or use no scaling; all rank $\geq$5.

\vspace{0.2em}
\footnotesize
Stage-1 MSE is not directly comparable across the $\Delta P^{0.25}$ vs.\ $\Delta P^{0.5}$ groups because
the normalised target $S/A_{\mathrm{scale}}$ differs; physical MAE (mm) is the correct cross-variant metric.
\end{frame}

% ──────────────────────────────────────────────────────────────
\begin{frame}{LONO Results: Per-Fold Breakdown}
\small
\begin{center}
\renewcommand{\arraystretch}{1.18}
\setlength{\tabcolsep}{5pt}
\begin{tabular}{@{}lrrrrr|r@{}}
\toprule
\textbf{Variant} &
\textbf{Nozzle0} & \textbf{Nozzle1} & \textbf{Nozzle4} & \textbf{Nozzle5} & \textbf{Nozzle2} &
\textbf{Mean} \\
& \footnotesize(1038) & \footnotesize(81) & \footnotesize(81) & \footnotesize(81) & \footnotesize(195) & \\
\midrule
\rowcolor{winrow}
\texttt{a\_dp050\_plus\_pressures} & 18.95 & 4.65 & 5.47 & 2.76 & 5.48 & \textbf{7.46} \\
\texttt{a\_dp025\_plus\_pressures} & 23.42 & 4.35 & 5.18 & 2.86 & 5.27 & 8.22 \\
\texttt{a\_dp050} & 18.74 & 3.32 & 5.37 & 7.93 & 5.98 & 8.27 \\
\texttt{a\_dp025} & 22.59 & 3.17 & 5.14 & 8.62 & 9.64 & 9.83 \\
\midrule
\rowcolor{warnrow}
\texttt{a\_dp050\_plus\_diameter} & 15.01 & 11.71 & 23.39 & 4.31 & 5.87 & 12.06 \\
\rowcolor{warnrow}
\texttt{a\_dp025\_plus\_diameter} & 20.90 & 14.16 & 20.67 & 4.74 & 6.60 & 13.41 \\
\rowcolor{warnrow}
\texttt{a\_dp050\_plus\_pressures\_diameter} & 22.34 & 11.79 & 26.14 & 3.23 & 7.80 & 14.26 \\
\rowcolor{warnrow}
\texttt{legacy\_9\_no\_scale} & 19.81 & 12.41 & 27.29 & 2.75 & 9.83 & 14.42 \\
\rowcolor{warnrow}
\texttt{a\_dp025\_plus\_pressures\_diameter} & 22.32 & 11.68 & 27.39 & 2.34 & 8.67 & 14.48 \\
\bottomrule
\end{tabular}
\end{center}

\vspace{0.2em}
\footnotesize
All values in mm (physical MAE on held-out nozzle test split). Numbers in parentheses are
sample-group counts per fold.
\end{frame}

% ──────────────────────────────────────────────────────────────
\begin{frame}{Discussion: Four Consistent Patterns}
\small
\begin{block}{1. \good{$\Delta P^{0.5}$ beats $\Delta P^{0.25}$ everywhere}}
All five dp050 variants outperform their dp025 counterparts with the same residual features.
The data-supported transitional exponent $\alpha=0.5$ encodes more physical variance into $A_{\mathrm{scale}}$,
leaving less unexplained residual for the MLP.
\end{block}

\begin{block}{2. \good{Log-pressure residuals give a small but consistent gain}}
\texttt{a\_dp050\_plus\_pressures} $<$ \texttt{a\_dp050}: the pressure logs capture
non-trivial variation at large Nozzle0 ($18.95$ vs.\ $18.74$ mm; modest) but are decisive
on the sparser folds. Net gain: $\sim$0.8\,mm mean MAE.
\end{block}

\begin{block}{3. \warn{$d_n$ as a standalone input hurts (even with $A$-scale present)}}
Adding \texttt{diameter\_mm\_z} to any $A$-scale variant increases MAE by $4$--$6$\,mm on
average. The network uses the redundant $d_n$ channel to memorise nozzle identity, which
degrades out-of-nozzle generalisation. This \emph{directly confirms} the earlier decision
to remove $d_n$ from the production feature set.
\end{block}

\begin{block}{4. \warn{Legacy 9 without scaling is competitive only on the easiest folds}}
\texttt{legacy\_9\_no\_scale} ranks 8th overall; its Nozzle4 MAE ($27.29$\,mm) is the worst
single-fold result in the table.
\end{block}
\end{frame}

% ──────────────────────────────────────────────────────────────
\begin{frame}{Discussion: Nozzle0 is the Hardest Fold}
\small
\begin{columns}[T]
\column{0.55\textwidth}
Nozzle0 (\texttt{BC20220627\_HZ\_Nozzle0}) dominates the mean because it contributes
\textbf{1038 sample groups} --- about 5$\times$ more than any other fold.

\vspace{0.5em}
Its MAE ranges from \SI{15.0}{mm} (best: \texttt{a\_dp050\_plus\_diameter}) to \SI{23.4}{mm}
across all variants, substantially above the \SI{2.8}--\SI{5.5}{mm} range seen on the smaller
folds.

\vspace{0.5em}
Two reasons:
\begin{enumerate}
  \item Nozzle0 spans the \emph{widest} operating-condition envelope --- holding it out removes
        the training data that anchors extreme pressure/angle combinations.
  \item It is the only nozzle family with a substantially different geometric design
        (first-generation vs.\ newer Nozzle1--8), so its OOD gap is inherently larger.
\end{enumerate}

\column{0.42\textwidth}
\begin{center}
\renewcommand{\arraystretch}{1.2}
\small
\begin{tabular}{@{}lrr@{}}
\toprule
\textbf{Fold} & \textbf{Groups} & \textbf{Best MAE} \\
\midrule
Nozzle0 & 1038 & 15.01\,mm \\
Nozzle2 & 195  & 5.27\,mm \\
Nozzle1 & 81   & 3.17\,mm \\
Nozzle4 & 81   & 5.14\,mm \\
Nozzle5 & 81   & 2.34\,mm \\
\bottomrule
\end{tabular}
\end{center}
\vspace{0.5em}
\footnotesize
Best MAE = minimum across all 9 variants for that fold.
\end{columns}
\end{frame}

% ──────────────────────────────────────────────────────────────
\begin{frame}{Conclusion: Feature Engineering is Faithful and Useful}
\small
\begin{block}{What the ablation proves}
\begin{enumerate}
  \item \textbf{Physics-based amplitude scaling is load-bearing.}
        Even the simplest $A$-scale variant (\texttt{a\_dp025}) beats the 9-feature legacy
        baseline by $\sim$4.6\,mm mean MAE.
  \item \textbf{The $\Delta P^{0.5}$ exponent is empirically correct.}
        Across all four residual-feature levels, dp050 variants rank above their dp025
        counterparts without exception.
  \item \textbf{Removing $d_n$ as a standalone input is the right choice.}
        It is already encoded in $A_{\mathrm{scale}}$; re-introducing it creates a memorisation
        pathway that costs 4--6\,mm OOD MAE.
  \item \textbf{Production variant \texttt{a\_dp050\_plus\_pressures} is validated.}
        It achieves the lowest mean MAE (7.46\,mm) and the most consistent per-fold
        performance of all nine configurations tested.
\end{enumerate}
\end{block}

\vspace{0.3em}
\textbf{Experiment:} \texttt{run\_stage1\_lono\_ablation.py} --- 9 variants $\times$ 5 LONO
folds, seed 42, \texttt{train\_stage1\_mse.py} per cell.\\
\textbf{Output:} \texttt{MLP/runs\_mlp/stage1\_lono\_ablation\_20260520\_220318/}
\end{frame}

\end{document}
